\section{Proxy}
\label{sec:pat-proxy}

\problemstatement{Provide a surrogate or placeholder for another object to
control access to it -- for lazy initialisation, access control, caching,
logging, or standing in for a remote object.}

\begin{patternbox}[When You Reach for It]
The trigger is \emph{``wrap an object to control how and when it is
used''} without changing its interface: delay an expensive load until
needed (virtual proxy), check permissions (protection proxy), cache
results, or represent something across a network (remote proxy).
\end{patternbox}

\subsection*{Understanding the pattern}

A Proxy implements the \emph{same interface} as the real subject and holds
a reference to it. It intercepts each call, does its extra job (lazy
create, permission check, cache lookup, log), then forwards to the real
object. Because the interface matches, clients cannot tell they are
talking to a proxy rather than the real thing.

\begin{ideabox}[Same Interface, Controlled Access]
Proxy shares Decorator's ``wrap with the same interface'' shape but its
intent is \emph{access control and management}, not adding features. The
canonical use is the virtual proxy: construct the heavy real object only
on first genuine use.
\end{ideabox}

\subsection*{C++ implementation}

\begin{lstlisting}
#include <bits/stdc++.h>
using namespace std;

struct Image {
    virtual void display() = 0;
    virtual ~Image() = default;
};

class RealImage : public Image {                // expensive real subject
    string filename;
public:
    RealImage(const string& f) : filename(f) {
        cout << "Loading " << filename << " from disk\n";   // heavy
    }
    void display() override { cout << "Displaying " << filename << "\n"; }
};

class ImageProxy : public Image {               // virtual proxy
    string filename;
    unique_ptr<RealImage> real;                 // created lazily
public:
    ImageProxy(const string& f) : filename(f) {}
    void display() override {
        if (!real) real = make_unique<RealImage>(filename);  // load on demand
        real->display();
    }
};

int main() {
    ImageProxy img("photo.png");                // no disk load yet
    cout << "Proxy created, nothing loaded.\n";
    img.display();                              // loads now, then displays
    img.display();                              // reuses, no reload
    return 0;
}
\end{lstlisting}

\begin{complexitybox}[Trade-offs]
\textbf{Pros.} Adds access control, laziness, caching or remoting
transparently; the client interface is unchanged. \textbf{Cons.} Extra
indirection and latency on first access; a poorly designed proxy can hide
performance costs or introduce inconsistency between proxy and subject.
\end{complexitybox}

\begin{takeawaybox}
Proxy stands in for a real object behind an identical interface to control
access -- lazy loading, protection, caching, or remoting. It looks like
Decorator but the intent differs: Decorator \emph{adds behaviour}, Proxy
\emph{governs access}. Name the proxy variant (virtual, protection,
remote) when you use it.
\end{takeawaybox}
